% Isolated production-verification candidate; not live-book source.
% Date: 30 August 2026.

\begin{theorem}
\label{thm:p071-infinite-star-structure}
Let \(B\) be a set and let
\[
 p:B\times[-\infty,0]\longrightarrow X_B
\]
be the CLD star of \cite[\S6.7.3]{CLD25}: all points \((b,0)\)
are identified with the centre \(\eta\), and an open neighbourhood of
\(\eta\) contains the whole branch
\(X_b=p(\{b\}\times[-\infty,0])\) for all but finitely many \(b\).
For a function \(v\) on a subset of \(X_B\), write
\(v_b(t)=v(p(b,t))\).

Let \(U\) be an open neighbourhood of \(\eta\).  If
\(u\in\operatorname{Psh}_{X_B}(U)\), then there is a countable set
\(C\subset B\) such that
\begin{equation}
\label{eq:p071-countable-full-branch-constancy}
 U\cap X_b=X_b,
 \qquad
 u_b(t)=u(\eta)
 \quad(t\in[-\infty,0])
\end{equation}
for every \(b\in B\setminus C\).  The same conclusion holds at every
ordinal stage \(\operatorname{Psh}_{X_B,\alpha}\) of
\cite[Lemma~8.1.4]{CLD25}, including when \(u(\eta)=-\infty\).

There is also a quantitative first-stage statement.  Let \(V\) be an open
neighbourhood of \(\eta\), put
\[
 A=\{a\in B:V\cap X_a\ne X_a\},
\]
and choose \(\rho_a>0\) such that
\([-\rho_a,0]\subset V_a\) for every \(a\in A\).  Thus \(A\) is finite.
Suppose that
\[
 f_n\in\operatorname{Psh}_{X_B,0}(V),
 \qquad f_{n+1}\leq f_n,
 \qquad f_n\longrightarrow u
\]
pointwise on \(V\), where
\(\operatorname{Psh}_{X_B,0}\) denotes the smooth functions with
\(\mathrm d'\mathrm d''f_n\geq0\).  Assume first that
\(c=u(\eta)\in\mathbf R\).  Then every \(u_a(-\rho_a)\), \(a\in A\), is
finite, and
\begin{equation}
\label{eq:p071-exceptional-slope-budget}
 R_A(u;\rho)
 :=\sum_{a\in A}
       \frac{u_a(-\rho_a)-c}{\rho_a}
 \in[0,+\infty).
\end{equation}
For every finite set \(J\subset B\setminus A\) and every choice of
\(r_b>0\), \(b\in J\), one has
\begin{equation}
\label{eq:p071-finite-secant-budget}
 0\leq
 \sum_{b\in J}\frac{c-u_b(-r_b)}{r_b}
 \leq R_A(u;\rho).
\end{equation}
For each \(b\in B\setminus A\), the function \(u_b\) is real-valued,
convex, and nondecreasing on the finite stratum
\((-\infty,0]\subset\mathbf R\), and
\begin{equation}
\label{eq:p071-two-endpoint-traces}
 \lim_{r\downarrow0}u_b(-r)=c,
 \qquad
 u_b(-\infty)=\lim_{r\to+\infty}u_b(-r)
 \quad\text{in }\mathbf R\cup\{-\infty\}.
\end{equation}
Its inward endpoint slope is therefore well defined by
\begin{equation}
\label{eq:p071-inward-endpoint-slope}
 \sigma_b(u)
 :=\lim_{r\downarrow0}\frac{c-u_b(-r)}{r}
 =\sup_{r>0}\frac{c-u_b(-r)}{r}
 \in[0,+\infty).
\end{equation}
If an arbitrarily indexed sum of nonnegative numbers is interpreted as the
supremum of its finite subsums, then
\begin{equation}
\label{eq:p071-summable-endpoint-slopes}
 \sum_{b\in B\setminus A}\sigma_b(u)
 \leq R_A(u;\rho)<+\infty.
\end{equation}
Consequently, \(u_b\) is nonconstant on at most countably many full
branches.

If instead \(u(\eta)=-\infty\), then \(u\) is identically \(-\infty\) on
the component of \(V\cap X_b\) containing \(\eta\), for every \(b\in B\).
No assertion is made about a disconnected component of \(V\cap X_b\).

The countable bound cannot be replaced by a finite bound, and the constant
in \eqref{eq:p071-summable-endpoint-slopes} is sharp.  More precisely, if
\(B\) contains distinct elements
\(a,b_1,b_2,\ldots\), there are a centre neighbourhood \(V\) with
\(V_a=(-1,0]\), smooth psh functions \(f_n\downarrow u\) on \(V\), and
\(u(\eta)=0\), such that every \(u_{b_j}\) is nonconstant and, for every
\(0<\rho<1\),
\begin{equation}
\label{eq:p071-sharp-slope-budget}
 \sum_{j\geq1}\sigma_{b_j}(u)
 =1
 =\frac{u_a(-\rho)-u(\eta)}{\rho}.
\end{equation}
\end{theorem}
